\documentclass[12pt, a4paper]{article}

% ============================================================
%  PACKAGES  (minimal set — compiles fast on Overleaf free)
% ============================================================
\usepackage[utf8]{inputenc}
\usepackage[T1]{fontenc}
\usepackage[vietnamese]{babel}
\usepackage{lmodern}
\usepackage{amsmath, amsthm, amssymb, mathtools}
\usepackage[margin=2.5cm]{geometry}
\usepackage[colorlinks=true, linkcolor=blue!70!black,
            citecolor=teal, urlcolor=blue]{hyperref}
\usepackage{enumitem}
\usepackage{booktabs}
\usepackage{xcolor}
\usepackage[most]{tcolorbox}
\usepackage{tikz}
\usetikzlibrary{arrows.meta, positioning}
\usepackage{algorithm}
\usepackage{algpseudocode}
\usepackage{microtype}

% ============================================================
%  COLOURS
% ============================================================
\definecolor{deepblue} {RGB}{10,  60, 120}
\definecolor{midblue}  {RGB}{30, 100, 180}
\definecolor{softred}  {RGB}{180, 40,  40}
\definecolor{softgreen}{RGB}{30, 120,  60}
\definecolor{amber}    {RGB}{180, 110,  10}
\definecolor{lightgray}{RGB}{245, 245, 248}

% ============================================================
%  TCOLORBOX THEOREM STYLES
% ============================================================
\tcbset{mybase/.style={
  breakable, enhanced jigsaw,
  fonttitle=\bfseries\small,
  left=4pt, right=4pt, top=3pt, bottom=3pt,
  before skip=8pt, after skip=8pt,
}}

\newtcbtheorem[number within=section]{theorem}{Định lý}{
  mybase, colback=deepblue!5, colframe=deepblue!60,
  attach boxed title to top left={xshift=4mm,yshift=-2mm},
  boxed title style={colback=deepblue!75, colframe=deepblue!75,
                     fontupper=\color{white}\bfseries\small}
}{thm}

\newtcbtheorem[number within=section]{lemma}{Bổ đề}{
  mybase, colback=midblue!5, colframe=midblue!50,
  attach boxed title to top left={xshift=4mm,yshift=-2mm},
  boxed title style={colback=midblue!65, colframe=midblue!65,
                     fontupper=\color{white}\bfseries\small}
}{lem}

\newtcbtheorem[number within=section]{definition}{Định nghĩa}{
  mybase, colback=softgreen!5, colframe=softgreen!55,
  attach boxed title to top left={xshift=4mm,yshift=-2mm},
  boxed title style={colback=softgreen!65, colframe=softgreen!65,
                     fontupper=\color{white}\bfseries\small}
}{def}

\newtcbtheorem[number within=section]{proposition}{Mệnh đề}{
  mybase, colback=amber!5, colframe=amber!55,
  attach boxed title to top left={xshift=4mm,yshift=-2mm},
  boxed title style={colback=amber!65, colframe=amber!65,
                     fontupper=\color{white}\bfseries\small}
}{prop}

\newtcbtheorem[number within=section]{remark}{Nhận xét}{
  mybase, colback=lightgray, colframe=gray!40,
  attach boxed title to top left={xshift=4mm,yshift=-2mm},
  boxed title style={colback=gray!35, colframe=gray!35,
                     fontupper=\bfseries\small}
}{rem}

\newtcolorbox{keyidea}{
  mybase, colback=softred!5, colframe=softred!65,
  title={\color{white}\bfseries Bài Học Cốt Lõi},
  attach boxed title to top left={xshift=4mm,yshift=-2mm},
  boxed title style={colback=softred!75, colframe=softred!75}
}

% ============================================================
%  MATH MACROS
% ============================================================
\newcommand{\HH}{\mathcal{H}}
\newcommand{\prox}{\operatorname{prox}}
\newcommand{\dom}{\operatorname{dom}}
\newcommand{\dist}{\operatorname{dist}}
\newcommand{\norm}[1]{\left\|#1\right\|}
\newcommand{\inner}[2]{\langle #1,\,#2\rangle}
\newcommand{\R}{\mathbb{R}}
\newcommand{\N}{\mathbb{N}}
\DeclareMathOperator*{\argmin}{arg\,min}

% ============================================================
%  TITLE
% ============================================================
\title{%
  \textcolor{deepblue}{\LARGE\bfseries
    Các Thuật Toán Điểm Kề Không Chính Xác Và Được Gia Tốc}\\[6pt]
  \large Một Bài Trình Bày Có Cấu Trúc Dành Cho Học Viên Cao Học\\[4pt]
  \normalsize
  Dựa trên: Salzo \& Villa (2012),
  \textit{J.\ Convex Anal.}\ \textbf{19}(4):1167--1192
}
\author{Ghi chú bài giảng của [Tên Của Bạn]}
\date{\today}

% ============================================================
\begin{document}

\begin{titlepage}
\centering
\vspace*{2cm}
{\huge\bfseries Inexact and Accelerated
Proximal Point Algorithms\par}
\vspace{2cm}
{\Large\bfseries Đề tài 10\par}
\vspace{0.5cm}
\begin{tabular}{ll}
Nguyễn Hoàng Tú & -- 23110220 \\
Bùi Công Hoàng Vũ & -- 23110223 \\
Huỳnh Trung Kiên & -- 22110091 \\
Nguyễn Ngọc Diễm Quỳnh & -- 21110167 \\
\end{tabular}\par
\vspace{2cm}
{\large\bfseries Giảng Viên Hướng Dẫn\par}
\vspace{0.2cm}
{\large Nguyễn Đăng Khoa\par}
\vfill
{\large Ngày 21 tháng 7 năm 2026\par}
\end{titlepage}

\maketitle\thispagestyle{empty}
\tableofcontents
\bigskip\hrule\bigskip

% ============================================================
\section*{Lời nói đầu: Tại sao bài báo này quan trọng}
\addcontentsline{toc}{section}{Lời nói đầu}
% ============================================================

Thuật toán điểm kề (Proximal Point Algorithms - PPA) là nền tảng của tối ưu hóa lồi hiện đại trong phục hồi ảnh, học máy và các bài toán ngược.
PPA \emph{gia tốc} của G\"{u}ler (1992) đạt được tốc độ hội tụ bậc nhất tối ưu
$F(x_k)-F^*=O(1/k^2)$ --- nhưng trong thực tế bài toán con điểm kề chỉ được giải một cách \emph{xấp xỉ}.
Salzo \& Villa (2012) mang đến ba đóng góp mang tính quyết định:

\begin{enumerate}[label=\textbf{C\arabic*.}, leftmargin=2.5em]
  \item \textbf{Sửa lại một lỗi nhỏ} trong phân tích không chính xác của G\"{u}ler, lỗi đã làm mất đi tính đúng đắn của khẳng định đạt tốc độ $O(1/k^2)$ dưới các tính toán không chính xác.
  \item Chỉ ra rằng \textbf{sai số loại 1} giới hạn tốc độ ở mức $O(1/k)$,
        bất kể chuỗi sai số tiến về $0$ nhanh đến mức nào.
  \item Giới thiệu \textbf{sai số loại 2}, nhờ đó khôi phục hoàn toàn tốc độ
        $O(1/k^2)$, ngay cả với một dãy sai số không khả tổng.
\end{enumerate}

\begin{keyidea}
\textbf{Bài học cốt lõi:}
\emph{Loại sai số cũng quan trọng không kém độ lớn của nó.}
Việc chuyển từ sai số loại 1 sang loại 2 giúp mở khóa tốc độ hội tụ $O(1/k^2)$.
Dưới sai số loại 1, việc gia tốc không mang lại lợi ích nào so với thuật toán PPA cổ điển.
\end{keyidea}

% ============================================================
\section{Thiết lập bài toán}
% ============================================================

\subsection{Bài toán Tối ưu hóa}

Xuyên suốt bài, $\HH$ là một không gian Hilbert thực với tích vô hướng $\inner{\cdot}{\cdot}$
và chuẩn $\norm{\cdot}$.
Chúng ta xét bài toán
\begin{equation}\label{eq:P}
  \min_{x\in\HH} F(x),
\end{equation}
trong đó $F:\HH\to\R\cup\{+\infty\}$ là hàm \emph{thực sự (proper), lồi,
và nửa liên tục dưới}.
Ta ký hiệu $F^*=\inf_x F(x)$ (có thể không đạt được cực trị).

\subsection{Toán tử Kề (Proximal Operator)}

\begin{definition}{Toán tử Kề}{prox}
Với $\lambda>0$ và $y\in\HH$, \emph{điểm kề} của $y$ đối với $\lambda F$ là
\[
  \prox_{\lambda F}(y)
  := \argmin_{x\in\HH}\!\left\{F(x)+\tfrac{1}{2\lambda}\norm{x-y}^2\right\}.
\]
\end{definition}

Đặt $\Phi_\lambda(x):=F(x)+\frac{1}{2\lambda}\norm{x-y}^2$,
điều kiện tối ưu bậc nhất cho ta:
\begin{equation}\label{eq:prox-equiv}
  z=\prox_{\lambda F}(y)
  \;\Longleftrightarrow\;
  \frac{y-z}{\lambda}\in\partial F(z)
  \;\Longleftrightarrow\;
  z=(I+\lambda\partial F)^{-1}(y).
\end{equation}

\paragraph{Tại sao lại có sự không chính xác?}
Trong các ứng dụng (khử mờ bằng tổng biến thiên, tính thưa có cấu trúc),
$\prox_{\lambda F}$ không có dạng biểu thức đóng (closed form) và phải được tính toán lặp,
do đó chỉ có thể tính \emph{xấp xỉ}. Việc hiểu rõ tiêu chí sai số nào bảo toàn được tốc độ hội tụ là rất quan trọng.

% ============================================================
\section{Ba Khái Niệm Về Sự Không Chính Xác}
% ============================================================

Cả ba tiêu chí đều được xây dựng dựa trên \emph{$\varepsilon$-dưới vi phân}.

\begin{definition}{$\varepsilon$-Dưới vi phân}{epsub}
Với $\varepsilon\ge 0$ và $z\in\dom F$:
\[
  \partial_\varepsilon F(z)
  :=\bigl\{\xi\in\HH : F(x)\ge F(z)+\inner{x-z}{\xi}-\varepsilon,
  \;\forall x\in\HH\bigr\}.
\]
Đặc biệt, $0\in\partial_\varepsilon F(z)\;\Leftrightarrow\;
F(z)\le F^*+\varepsilon$.
\end{definition}



\subsection{Xấp xỉ Loại 1}

\begin{definition}{Xấp xỉ Loại 1 (Rockafellar / G\"{u}ler)}{type1}
$z\approx_1\prox_{\lambda F}(y)$ với độ chính xác $\varepsilon$ khi và chỉ khi
\[
  0\in\partial_{\varepsilon^2/(2\lambda)}\Phi_\lambda(z).
  \tag{AT1}
\]
Tương đương với việc, $z$ đạt được giá trị gần cực tiểu của hàm lồi mạnh $\Phi_\lambda$
với độ gần tối ưu (suboptimality) $\le\varepsilon^2/(2\lambda)$, và $\norm{z-\prox_{\lambda F}(y)}\le\varepsilon$.
\end{definition}


\subsection{Xấp xỉ Loại 2}

\begin{definition}{Xấp xỉ Loại 2 (Burachik--Iusem--Svaiter)}{type2}
$z\approx_2\prox_{\lambda F}(y)$ với độ chính xác $\varepsilon$ khi và chỉ khi
\[
  \frac{y-z}{\lambda}\in\partial_{\varepsilon^2/(2\lambda)}F(z).
  \tag{AT2}
\]
Điều này nới lỏng trực tiếp quan hệ bao hàm trong \eqref{eq:prox-equiv};
tương đương với việc $z\in(I+\lambda\partial_{\varepsilon^2/(2\lambda)}F)^{-1}(y)$,
một sự mở rộng $\varepsilon$ của ánh xạ kề.
\end{definition}


\subsection{Xấp xỉ Loại 3}

\begin{definition}{Xấp xỉ Loại 3 (Rockafellar 1976)}{type3}
$z\approx_3\prox_{\lambda F}(y)$ với độ chính xác $\varepsilon$ khi và chỉ khi
\[
  \dist(0,\partial\Phi_\lambda(z))\le\varepsilon/\lambda,
  \tag{AT3}
\]
tương đương với việc, tồn tại $\exists\,e\in\HH$, $\norm{e}\le\varepsilon$, sao cho
$z=\prox_{\lambda F}(y+e)$.
Do đó, xấp xỉ loại 3 chính là điểm kề \emph{chính xác} của một đầu vào bị nhiễu.
\end{definition}



\subsection{Phân Cấp Giữa Các Loại Sai Số}

\begin{proposition}{Sơ đồ Quan hệ}{hier}
\begin{enumerate}[label=(\roman*)]
  \item $z\approx_2\prox_{\lambda F}(y)$ với độ chính xác $\varepsilon$
        $\Rightarrow z\approx_1\prox_{\lambda F}(y)$ với độ chính xác $\varepsilon$.
  \item $z\approx_3\prox_{\lambda F}(y)$ với độ chính xác $\varepsilon$
        $\Rightarrow z\approx_1\prox_{\lambda F}(y)$ với độ chính xác $\varepsilon$.
  \item Nếu $F$ lồi mạnh (với mô-đun $\mu>0$):
        loại 1 với độ chính xác $\varepsilon$
        $\Rightarrow$ loại 2 với độ chính xác $\varepsilon\sqrt{(1+\lambda\mu)/(\lambda\mu)}$.
\end{enumerate}
Do đó, loại 2 và loại 3 là các trường hợp đặc biệt khắt khe hơn của loại 1.
\end{proposition}

\begin{center}
\begin{tikzpicture}[
  node distance=1.6cm and 2.8cm,
  box/.style={draw, rounded corners=4pt, minimum width=3cm,
              minimum height=0.9cm, align=center, font=\small},
  arr/.style={-{Stealth[length=5pt]}, thick, shorten >=2pt, shorten <=2pt}
]
  \node[box, fill=softred!12,   draw=softred!70]  (t1)
        {Loại 1\\[-2pt]\tiny(ít khắt khe nhất)};
  \node[box, fill=midblue!10,   draw=midblue!60,
        right=of t1]             (t2) {Loại 2};
  \node[box, fill=softgreen!10, draw=softgreen!60,
        below=1.2cm of t1]       (t3) {Loại 3};
  \draw[arr, softred!60]   (t2) -- node[above,font=\tiny]{suy ra} (t1);
  \draw[arr, softred!60]   (t3) -- node[right,font=\tiny]{suy ra} (t1);
  \draw[arr, dashed, gray] (t1)
        to[bend left=15] node[above,font=\tiny,sloped]{lồi mạnh} (t2);
\end{tikzpicture}
\end{center}

% ============================================================
\section{Dãy Ước Lượng của Nesterov}
% ============================================================

\begin{definition}{Dãy Ước lượng}{estseq}
$\bigl((\varphi_k),(\beta_k)\bigr)$ là một \emph{dãy ước lượng} của $F$ khi và chỉ khi
\[
  \varphi_k(x)-F(x)\le\beta_k(\varphi_0(x)-F(x)),\quad
  \forall x\in\HH,\;\forall k,\quad
  \text{và}\quad\beta_k\to 0.
\]
\end{definition}


\begin{theorem}{Sự hội tụ thông qua Dãy ước lượng}{estconv}
Giả sử $\bigl((\varphi_k),(\beta_k)\bigr)$ là một dãy ước lượng của $F$.
Nếu $(x_k)$ và $(\delta_k\ge 0)$ thỏa mãn
$F(x_k)\le(\varphi_k)^*+\delta_k$,
thì với mọi $x\in\dom F$:
\[
  F(x_k)\le\beta_k(\varphi_0(x)-F(x))+\delta_k+F(x).
\]
Nếu $\delta_k\to 0$, thì $(x_k)$ là một dãy cực tiểu hóa.
Hơn nữa, nếu $F^*=F(x^*)$ tại một $x^*$ nào đó:
\[
  \boxed{F(x_k)-F^*\le\beta_k(\varphi_0(x^*)-F^*)+\delta_k.}
\]
\end{theorem}

\subsection{Toán Tử Cập Nhật Và Tính Đóng Bậc Hai}

Định nghĩa với $\xi\in\partial_\varepsilon F(z)$ và $\alpha\in[0,1)$:
\[
  \mathcal{U}(z,\varepsilon,\xi,\alpha)(\varphi)(x)
  :=(1-\alpha)\varphi(x)+\alpha\bigl(F(z)+\inner{x-z}{\xi}-\varepsilon\bigr).
\]
Toán tử này thỏa mãn $\widehat\varphi(x)-F(x)\le(1-\alpha)(\varphi(x)-F(x))$,
đây chính xác là bất đẳng thức đệ quy của dãy ước lượng.

\textbf{Tính đóng bậc hai.}
Nếu $\varphi(x)=\varphi^*+\frac{A}{2}\norm{x-\nu}^2$, thì
$\mathcal{U}(z,\varepsilon,\xi,\alpha)(\varphi)$ lại là một hàm bậc hai với:
\begin{equation}\label{eq:quadupdate}
  \hat\varphi^*=(1-\alpha)\varphi^*+\alpha F(z)+\alpha\inner{\nu-z}{\xi}
               -\tfrac{\alpha^2}{2(1-\alpha)A}\norm{\xi}^2-\alpha\varepsilon,
  \quad
  \hat A=(1-\alpha)A,\quad
  \hat\nu=\nu-\tfrac{\alpha}{(1-\alpha)A}\xi.
\end{equation}

\begin{lemma}{Tốc độ Giảm của $\beta_k$}{betarate}
Cho $\lambda_k>0$, $A>0$, $0<a\le b$.
Định nghĩa $A_{k+1}=(1-\alpha_k)A_k$ và
$\alpha_k^2/\bigl((1-\alpha_k)A_k\lambda_k\bigr)\in[a,b]$.
Khi đó $\beta_k:=\prod_{i=0}^{k-1}(1-\alpha_i)$ thỏa mãn:
\[
  \frac{1}{\!\left(1+\sqrt{bA}\sum_{j=0}^{k-1}\sqrt{\lambda_j}\right)^{\!2}}
  \le\beta_k\le
  \frac{1}{\!\left(1+\frac{\sqrt{aA}}{2}\sum_{j=0}^{k-1}\sqrt{\lambda_j}\right)^{\!2}}.
\]
Đặc biệt, $\lambda_k\ge\lambda>0$ với mọi $k$ suy ra $\beta_k=O(1/k^2)$.
\end{lemma}

% ============================================================
\section{Phân Tích Hội Tụ}
% ============================================================

\subsection{Thuật toán IAPPA1: Sai số Loại 1}

Khó khăn cốt lõi: bước cập nhật của G\"{u}ler
$\hat\nu=\nu-(t-z)/\alpha$ sử dụng $z=\prox_{\lambda F}(t)$,
giá trị này là \emph{không xác định được}.
Cách khắc phục là theo dõi một biến phụ $u_k\approx\nu_k$ với
nhiễu được định lượng $\eta_k=\norm{u_k-\nu_k}$.

\begin{algorithm}[H]
\caption{IAPPA1 --- Thuật toán PPA Gia tốc Không chính xác với Sai số Loại 1}
\label{alg:iappa1}
\begin{algorithmic}[1]
\Require $x_0\in\dom F$, $u_0=x_0$, $A_0=A>0$,
         $\eta_0=\delta_0=0$, các dãy $(\lambda_k),(\varepsilon_k)>0$.
\For{$k=0,1,2,\ldots$}
  \State $\alpha_k\leftarrow$ nghiệm dương của $\alpha^2=(1-\alpha)A_k\lambda_k$
  \State $y_k\leftarrow(1-\alpha_k)x_k+\alpha_k u_k$
  \State $x_{k+1}\approx_1\prox_{\lambda_k F}(y_k)$
         với độ chính xác $\varepsilon_k$
  \State $A_{k+1}\leftarrow(1-\alpha_k)A_k$
  \State $u_{k+1}\leftarrow u_k-\tfrac{1}{\alpha_k}(y_k-x_{k+1})$
  \State $\eta_{k+1}\leftarrow\eta_k+\varepsilon_k/\alpha_k$
  \State $\delta_{k+1}\leftarrow(1-\alpha_k)\delta_k
         +\tfrac{(\alpha_k\eta_{k+1})^2}{2\lambda_k}$
\EndFor
\end{algorithmic}
\end{algorithm}

\begin{theorem}{Sự Hội tụ của IAPPA1}{iappa1}
Giả sử $\lambda_j\le M\lambda_i$ với mọi $j\le i$.
Nếu $\varepsilon_k=O(1/k^q)$ với $q>3/2$, thì $(x_k)$ là một dãy cực tiểu hóa.
Nếu đạt được giá trị $F^*$:
\[
  F(x_k)-F^* =
  \begin{cases}
    O(1/k^{2q-3}) & 3/2<q<2,\\
    O(\log^2\!k/k) & q=2,\\
    O(1/k)         & q>2.
  \end{cases}
\]
\end{theorem}

\begin{remark}{Tại sao gia tốc thất bại dưới sai số loại 1}{fail}
Từ $\eta_k=\sum_{i<k}\varepsilon_i/\alpha_i$ và $1/\alpha_i=O(i)$,
ngay cả khi $\varepsilon_i=O(1/i^2)$ ta cũng có $\eta_k=O(\log k)$,
và $\delta_k=O(\log^2\!k/k)$.
Tốc độ $O(1/k^2)$ là không thể đạt được dưới sai số loại 1,
bất kể $\varepsilon_k\to 0$ nhanh như thế nào.
Việc gia tốc \emph{không mang lại lợi thế nào} so với PPA cơ bản trong trường hợp này.
\end{remark}

\subsection{Thuật toán IAPPA2: Sai số Loại 2}

Với sai số loại 2, $\xi_{k+1}=(y_k-x_{k+1})/\lambda_k$ có thể được tính toán trực tiếp
và $\xi_{k+1}\in\partial_{\varepsilon_k^2/(2\lambda_k)}F(x_{k+1})$.
Bước cập nhật cho $\delta_k$ được đơn giản hóa đi rất nhiều.

\begin{algorithm}[H]
\caption{IAPPA2 --- Thuật toán PPA Gia tốc Không chính xác với Sai số Loại 2}
\label{alg:iappa2}
\begin{algorithmic}[1]
\Require $x_0=\nu_0\in\dom F$, $A_0=A>0$, $\delta_0=0$, $a\in(0,2]$,
         các dãy $(\lambda_k),(\varepsilon_k)>0$.
\For{$k=0,1,2,\ldots$}
  \State Chọn $\alpha_k\in[0,1)$ sao cho
         $a\le\tfrac{\alpha_k^2}{(1-\alpha_k)A_k\lambda_k}\le 2$
  \State $y_k\leftarrow(1-\alpha_k)x_k+\alpha_k\nu_k$
  \State $x_{k+1}\approx_2\prox_{\lambda_k F}(y_k)$
         với độ chính xác $\varepsilon_k$
  \State $A_{k+1}\leftarrow(1-\alpha_k)A_k$
  \State $\nu_{k+1}\leftarrow\nu_k
         -\tfrac{\alpha_k}{(1-\alpha_k)A_k\lambda_k}(y_k-x_{k+1})$
  \State $\delta_{k+1}\leftarrow(1-\alpha_k)\delta_k
         +\tfrac{\varepsilon_k^2}{2\lambda_k}$
\EndFor
\end{algorithmic}
\end{algorithm}

Dạng biểu thức đóng của sự tích lũy sai số là:
\begin{equation}\label{eq:delta2}
  \delta_k = \frac{\beta_k}{2}
             \sum_{i=0}^{k-1}\frac{\varepsilon_i^2}{\lambda_i\beta_{i+1}}.
\end{equation}

\begin{theorem}{Sự Hội tụ của IAPPA2}{iappa2}
Giả sử $\lambda_j\le M\lambda_i$ với mọi $j\le i$.
Nếu $\varepsilon_k=O(1/k^q)$ với $q>1/2$, thì $(x_k)$ là một dãy cực tiểu hóa.
Nếu đạt được giá trị $F^*$:
\[
  F(x_k)-F^* =
  \begin{cases}
    O(1/k^2)                       & q>3/2,\\
    O(1/k^2)+O(\log k/k^2)         & q=3/2,\\
    O(1/k^2)+O(1/k^{2q-1})         & 1/2<q<3/2.
  \end{cases}
\]
Cụ thể, $q>3/2$ khôi phục hoàn toàn tốc độ tối ưu $O(1/k^2)$.
\end{theorem}

\begin{remark}{Sai số không khả tổng}{nonsumm}
Với $1/2<q\le 1$, ta có $\sum_k\varepsilon_k=+\infty$,
nhưng sự hội tụ vẫn được đảm bảo.
Điều này hoàn toàn trái ngược với PPA không gia tốc (Rockafellar 1976),
vốn yêu cầu $\sum_k\varepsilon_k<+\infty$.
\end{remark}

\subsection{Tóm tắt: So sánh Tốc độ Hội tụ}

\begin{center}
\renewcommand{\arraystretch}{1.35}
\begin{tabular}{@{}llll@{}}
\toprule
\textbf{Thuật toán} & \textbf{Loại Sai số} & \textbf{Mức độ giảm} & \textbf{Tốc độ} \\
\midrule
PPA gia tốc chính xác & ---     & ---              & $O(1/k^2)$      \\
\midrule
IAPPA1              & Loại 1  & $O(k^{-q})$, $q>2$         & $O(1/k)$        \\
IAPPA1              & Loại 1  & $O(k^{-q})$, $3/2<q\le 2$ & $O(1/k^{2q-3})$ \\
\midrule
IAPPA2              & Loại 2  & $O(k^{-q})$, $q>3/2$          & $O(1/k^2)$      \\
IAPPA2              & Loại 2  & $O(k^{-q})$, $1/2<q\le 3/2$  & $O(1/k^{2q-1})$ \\
\bottomrule
\end{tabular}
\end{center}

% ============================================================
\section{Các Dạng Đơn Giản Tương Đương}
% ============================================================

Thay thế $t_k:=1/\alpha_k$, IAPPA2 trở thành:
\begin{equation}\label{eq:simple}
\left\{\;
\begin{aligned}
  t_{k+1} &= \tfrac{1+\sqrt{1+4a_k\lambda_k t_k^2/(a_{k+1}\lambda_{k+1})}}{2},\\[4pt]
  x_{k+1} &\approx_2\prox_{\lambda_k F}(y_k),\\[4pt]
  y_{k+1} &= x_{k+1}+\tfrac{t_k-1}{t_{k+1}}(x_{k+1}-x_k)
             +(1-a_k)\tfrac{t_k}{t_{k+1}}(y_k-x_{k+1}),
\end{aligned}
\right.
\end{equation}
trong đó $(a_k)\subset(0,2]$ được chọn tự do.
Đặt $a_k=2$, $\lambda_k=\lambda$ (hằng số), và sử dụng các bước kề chính xác sẽ khôi phục lại thuật toán \textbf{FISTA} của Beck \& Teboulle (2009).

% ============================================================
\section{Lỗi Trong Chứng Minh Của G\"{u}ler}
% ============================================================

\begin{remark}{Chứng minh của G\"{u}ler sai ở đâu}{guler}
Định lý 3.1 của G\"{u}ler (1992) khẳng định đạt được tốc độ $O(1/k^2)$ dưới sai số loại 1.
Lỗ hổng: bước cập nhật
$\hat\nu=\nu-(t-z)/\alpha$
sử dụng $z=\prox_{\lambda F}(t)$, giá trị này là \emph{không xác định được}.

Thay vào đó, việc dùng $x\approx_1\prox_{\lambda F}(t)$ đưa vào một lượng sai số
$\Delta=(x-z)/\alpha$, và lượng này tích lũy lại:
\[
  \eta_k=\sum_{i<k}\frac{\varepsilon_i}{\alpha_i}=O(k^{1-q}\text{ hoặc }\log k).
\]
Vì $\delta_k=O(\beta_k\sum_i\eta_{i+1}^2)$, điều này không thể là $O(1/k^2)$
với bất kỳ cách chọn $q$ nào. Salzo \& Villa đã sửa lỗi này bằng cách theo dõi rõ ràng
$\eta_k$ và thiết lập Định lý~\ref{thm:iappa1}.
\end{remark}

% ============================================================
\section{Tổng Kết và Các Câu Hỏi Mở}
% ============================================================

\paragraph{Những điểm chính rút ra.}
\begin{enumerate}[label=\textbf{T\arabic*.}, leftmargin=2.5em]
  \item Sai số loại 1 giới hạn tốc độ hội tụ ở mức $O(1/k)$; việc gia tốc là vô ích.
  \item Sai số loại 2 với $\varepsilon_k=O(1/k^q)$, $q>3/2$, cho tốc độ $O(1/k^2)$.
  \item Sai số loại 2 cho phép dãy $(\varepsilon_k)$ không khả tổng mà vẫn hội tụ.
  \item Đây là bài báo đầu tiên có phân tích hội tụ \emph{chính xác} cho thuật toán PPA gia tốc không chính xác.
\end{enumerate}

\paragraph{Các câu hỏi mở.}
\begin{itemize}
  \item \emph{Sự hội tụ của dãy lặp:} Bản thân dãy $(x_k)$ có hội tụ không (không chỉ $F(x_k)$)?
  \item \emph{Sai số tương đối:} Có thể mở rộng phân tích cho các tiêu chí xấp xỉ tương đối không?
  \item \emph{Tính lồi mạnh:} Tốc độ hội tụ tuyến tính $O(\rho^k)$ có được bảo toàn dưới sự không chính xác không?
  \item \emph{Sai số ngẫu nhiên:} Các mở rộng cho xấp xỉ ngẫu nhiên (MCMC / dưới đạo hàm ngẫu nhiên)?
\end{itemize}

% ============================================================
\begin{thebibliography}{9}

\bibitem{sv2012}
S.~Salzo, S.~Villa,
\textit{Inexact and Accelerated Proximal Point Algorithms},
J.\ Convex Anal.\ \textbf{19}(4) (2012) 1167--1192.

\bibitem{guler1992}
O.~G\"{u}ler,
\textit{New proximal point algorithms for convex minimization},
SIAM J.\ Optim.\ \textbf{2}(4) (1992) 649--664.

\bibitem{nesterov2004}
Y.~Nesterov,
\textit{Introductory Lectures on Convex Optimization},
Kluwer, Boston (2004).

\bibitem{rockafellar1976}
R.~T.~Rockafellar,
\textit{Monotone operators and the proximal point algorithm},
SIAM J.\ Control Optim.\ \textbf{14}(5) (1976) 877--898.

\bibitem{bt2009}
A.~Beck, M.~Teboulle,
\textit{A fast iterative shrinkage-thresholding algorithm},
SIAM J.\ Imaging Sci.\ \textbf{2}(1) (2009) 183--202.

\end{thebibliography}

\end{document}
